\section{Low degrees and comparison with other Gersten statements}\label{d3:sec:comparisons}
We use the arithmetic-quadric notation of Section~\ref{d3:sec:geometry} in this appendix.

\subsection{A low-degree stress test}
The analogous $K_1$ construction does not supply a counterexample.
Its product regulator would have degree three and twist two.
Here the filtration is the analytic truncation of the holomorphic
de Rham complex, which need not vanish in filtration degree two
on a complex surface.

For a smooth affine complex surface $U$, write $X=U^{\mathrm{an}}$.
An affine embedding realizes $X$ as a closed complex submanifold
of $\C^N$. Coordinate functions separate points and give local
coordinates. They bound the holomorphic hull of each compact set
inside a compact coordinate polydisc; the hull is closed there
and hence compact. Together with second countability, these
properties make $X$ Stein.
Cartan's theorem~B
\cite[\S5, proof of Theorem~B, pp.~19-14--19-15]{CartanStein1952}
gives $H^q(X,\Omega_X^a)=0$ for every $q>0$.
The holomorphic Poincar\'e lemma follows on coordinate balls by
radial integration and identifies $\C_X$ with the holomorphic
de Rham complex. Its hypercohomology is therefore computed by
global holomorphic forms, in degrees zero through two. Thus
\[
 H^2(X,\C)=\Gamma(X,\Omega_X^2)/d\Gamma(X,\Omega_X^1),
 \qquad H^3(X,\C)=0.
\]
For
$F^2_{\mathrm{an}}\Omega_X^\bullet=\Omega_X^{\ge2}
=\Omega_X^2[-2]$, the map from degree-two hypercohomology to
$H^2(X,\C)$ is the displayed quotient map and is surjective.
Moreover,
$\mathbb H^3(X,F^2_{\mathrm{an}}\Omega_X^\bullet)
=H^1(X,\Omega_X^2)=0$.
The constant real sheaf $\R(2)$ is a direct summand of the
underlying real sheaf $\C$, so $H^3(X,\R(2))=0$ as well.
The analytic Deligne cone sequence contains
\[
\begin{split}
 H^2(X,\R(2))\oplus
 \mathbb H^2(X,F^2_{\mathrm{an}}\Omega_X^\bullet)
 &\longrightarrow H^2(X,\C)\longrightarrow\HD^3(X,\R(2))\\
 &\longrightarrow H^3(X,\R(2))\oplus
 \mathbb H^3(X,F^2_{\mathrm{an}}\Omega_X^\bullet).
\end{split}
\]
Its first arrow is surjective and its final group is zero.
Hence $\HD^3(U,\R(2))=0$.
This calculation uses the analytic truncation just defined;
it makes no corresponding assertion about the logarithmic
mixed-Hodge filtration.
The elementary quotient in Lemma~\ref{d3:lem:Deligne} is
unavailable in twist two. This agrees with the fact that $K_1$
of a commutative local ring is its unit group and injects into
the fraction field.

There is also an explicit algebraic version of this test. On the split quadric in coordinates
\[
uv=t(1-t),
\]
the matrix
\[
e=\begin{pmatrix}t&u\\v&1-t\end{pmatrix}
\]
is a rank-one idempotent. Up to choosing the ruling, its image is the line bundle $P=\OO(D_+)$. For $\alpha\in E^\times$, $\alpha\neq1$, the determinant-one matrix
\[
M=\begin{pmatrix}\alpha^{-1}&0\\0&1\end{pmatrix}
\bigl(I+(\alpha-1)e\bigr)
\]
represents $\alpha([P]-1)$ in $K_1$. Its upper-left entry is $s/\alpha$ with $s=1+(\alpha-1)t$. On $D(s)$ Gaussian elimination makes $M$ elementary. The removed curve has
\[
t=-\frac1{\alpha-1},\qquad uv=-\frac{\alpha}{(\alpha-1)^2}\neq0,
\]
so it is a geometrically irreducible principal curve isomorphic to $\mathbb G_m$. Hence ordinary Betti injectivity alone cannot protect a $K_1$ product. The condition $3>\dim_{\C}U=2$ in the main proof is essential.

\subsection{Equal characteristic and other localizations}
Replacing the arithmetic DVR by $k[t]_{(t)}$ and adjoining $u$
with $u^2=-t$ does not reproduce the same coefficient class.
For every field $k$, restriction gives an isomorphism
\[
 K_3(k)_{\Q}^{(2)}
       \xrightarrow{\sim}K_3(k(u))_{\Q}^{(2)}.
\]
To prove this, put $F=k(u)$ and compare the Milnor and Quillen
rational-function residue sequences
\cite[III, Theorem~7.4; V, Corollary~6.7.1]{Weibel}.
They have respective kernels $K_3^M(k)$ and $K_3(k)$ and
surjective residue maps to
$\bigoplus_P K_2^M(k[u]/P)$ and
$\bigoplus_P K_2(k[u]/P)$, where $P$ runs over monic
irreducible polynomials.
The right-hand terms agree by Matsumoto's theorem
\cite[III, Theorem~6.1]{Weibel}.

Their residue maps are compatible with the symbol maps.
At each polynomial valuation, a product of unit classes extends
over the valuation ring and has zero boundary, while for a
uniformizer $\pi$ and units $a,b$ the boundary-product formula gives
\[
 \partial\{\pi,a,b\}=\{\bar a,\bar b\}.
\]
Multilinearity and $\{\pi,\pi\}=\{-1,\pi\}$ reduce every Milnor
symbol to unit symbols and symbols with exactly one uniformizer.
These formulas therefore give its tame residue
\cite[III, Theorem~7.3; V, equation~(6.1.1),
Example~6.1.2, and Lemma~6.7.3]{Weibel}.

Given $z\in K_3(F)$, surjectivity of the Milnor residue map
supplies an image $m$ of $K_3^M(F)$ with the same residues as
$z$. Exactness of the Quillen sequence gives
$z=\operatorname{res}(\beta)+m$ for some $\beta\in K_3(k)$.
The rational Milnor image has pure Adams weight three
\cite[IV, Example~5.9.1]{Weibel}.
Naturality of the rational Adams decomposition
\cite[IV, Theorem~5.11]{Weibel} therefore makes projection
to weight two surjective from $K_3(k)_{\Q}^{(2)}$.
Injectivity follows from the injective constant-field map in
the Quillen sequence, proving the displayed isomorphism.
The involution $u\mapsto-u$ fixes $k$ and hence acts trivially
on this space; in characteristic two it is already the identity.
Thus the needed rational weight-two anti-invariant regulator
source is absent. The arithmetic number-field class is an
essential ingredient.

The scheme $T$ is regular and contains a field, but it is not semilocal: its nonzero Picard group already rules that out. Local Gersten theorems over fields therefore do not give an injection $K_3(T)\to K_3(\Frac T)$ for this global scheme.

Sections~\ref{d3:sec:node}--\ref{d3:sec:Delignedetection} prove persistence for the precise noncomplete localization; they do not by themselves justify completion. Section~\ref{d3:sec:completed} supplies a separate arithmetic $p$-adic proof for completed examples. Propositions~\ref{d3:prop:explicitcomplete}, \ref{d3:prop:explicittwocomplete}, and \ref{d3:prop:explicitoddcomplete} verify that the particular explicit classes at every residue prime survive both completion and henselization. These arguments use arithmetic detectors after completion.

\subsection{The distinction between the generic fiber and the fraction field}
The rings considered here are regular and flat over $\Z$, and are consequently $F$-smooth in the sense relevant to \cite{BouisMotivic}. Corollaries~5.12--5.13 of that reference apply to them: in degree $i$, the motivic comparison with $\Z/p^\nu(i)$-coefficients injects into the cohomology of the entire generic fiber $A[1/p]$. This is compatible with the construction here, whose nonzero classes are detected on that generic fiber and become zero only after a further restriction to the fraction field. The generic fiber has positive dimension and cannot be replaced by its generic point in that statement.

The special fiber $A/(p)$ of the basic examples is the singular domain $k[x,y,z]/(xy+z^2)$, or its corresponding completion. Thus these maps are not smooth over their coefficient DVR. They cannot become smooth through a different local DVR coefficient presentation either: primality of $(p)$ forces the base uniformizer to be associated to $p$, and the same singular special fiber remains. Absolute purity also does not assert that a principal-divisor Gysin map in higher $K$-theory is zero; triviality of the normal bundle gives the weaker self-intersection identity $i^*i_*=0$.

\subsection{Rational and finite coefficients}
The regulator proves that the kernel contains an element of infinite order. It does not by itself prove that any fixed finite-coefficient Gersten map fails. In the explicit $p=3$ case, Section~\ref{d3:sec:modthree} supplies a separate algebraic proof for $\Z/3$-coefficients. That coefficient characteristic equals the residue characteristic; no prime-to-residue-characteristic conclusion is asserted here. A rationally nonzero class can otherwise become divisible in a sufficiently large group. In particular, one may not interchange a filtered colimit over shrinking opens with an inverse limit over coefficient exponents without a separate argument.

Feld's Theorem~1.1 \cite{FeldFiniteGersten2026} constructs a nonzero
generic-kernel class in $K_2(A_{\mathrm F};\mathbb{Z}/3)$, where
$A_{\mathrm F}=V[[x,y]]/(3+x^2-y^3)$ and $V$ is a complete mixed-characteristic
$(0,3)$ discrete valuation ring with uniformizer $3$.
In that example, an integral lift acquires divisibility by three over
the fraction field while the integral $K_2$ restriction is injective
\cite[Corollary~5.1]{FeldFiniteGersten2026}.

\subsection{An explicit number-field source when the prime is three}
Let $u^2=-3$ and $\zeta=(1+u)/2$. Then $\zeta=e^{i\pi/3}$ at the chosen embedding and
\[
1-\zeta=\zeta^{-1}.
\]
Thus the pre-Bloch symbol $[\zeta]$ has zero wedge boundary. Its Bloch--Wigner dilogarithm is positive, since
\[
D(e^{i\pi/3})=-\int_0^{\pi/3}\log\bigl(2\sin(t/2)\bigr)\,dt>0.
\]
The rational Bloch-group description of $K_3^{\mathrm{ind}}(E)$ \cite[VI, Theorem~5.2]{Weibel} supplies an explicit possible rational source class. The main proof does not require fixing this comparison or its integral normalization: Borel's theorem and Lemma~\ref{d3:lem:beta} already produce an integral anti-invariant source with nonzero regulator.

\section{A mixed-square obstruction to a simplicial-functor construction}
\label{d3:sec:mochizukisquare}

Theorem~5 of \cite{MochizukiParameters} asserts that, for a noetherian
local ring $B$ and a non-zero-divisor $g$, the functor
$\mathcal P_{B/gB}\to\mathcal M_B(1)$ induces the zero map on
$K$-theory. Its application to $B=A$ and $g=x$ would kill the supported
class used above. We therefore examine a specific step in its proof.
The following calculation concerns Lemma~13(III) as written; it does
not, by itself, establish that Theorem~5 is false.

Let $B$ be a local ring, let $g$ be a non-zero-divisor in its maximal
ideal, and put $t=(1+g)^{-1}$. Write
\[
G=\operatorname{Typ}(g)=[B\xrightarrow{g}B],\qquad
I=\operatorname{Typ}(1)=[B\xrightarrow{1}B],\qquad M=G\oplus I,
\]
with complexes in homological degrees $1,0$. These are respectively
the objects $(1,0)_B$, $(0,1)_B$, and $(1,1)_B$ of the category $\mathcal C$
in \cite[Definition~9]{MochizukiParameters}. For $b\in B$ define
chain maps $\alpha_b:I\to M$ and $\beta_b:M\to G$ by
\begin{align*}
(\alpha_b)_1&=\binom{b}{1},&
(\alpha_b)_0&=\binom{gb}{1},\\
(\beta_b)_1&=\begin{pmatrix}1&-b\end{pmatrix},&
(\beta_b)_0&=\begin{pmatrix}1&-gb\end{pmatrix}.
\end{align*}
They give a split short exact sequence of complexes
\[
\mathcal S_b:\qquad 0\longrightarrow I\xrightarrow{\alpha_b}
M\xrightarrow{\beta_b}G\longrightarrow0.
\]
Indeed $s:G\to M$, given by $\binom10$ in both degrees, is a
chain section of $\beta_b$, while $r:M\to I$, given by
$\begin{pmatrix}0&1\end{pmatrix}$ in both degrees, satisfies
$r\alpha_b=1$ and $\alpha_b r+s\beta_b=1$. In the block notation
of \cite[Conventions~11]{MochizukiParameters}, the arrows are
$\binom b1$ and $\begin{pmatrix}1&-b\end{pmatrix}$, respectively.
Both are upper triangular. Thus $\mathcal S_b$ is an object of
$S_2^{\triangle}\mathcal C$ in that paper's notation.

There is an isomorphism of short exact sequences
$\mathcal S_1\to\mathcal S_t$ whose components on $I,M,G$ are,
respectively,
\[
\ell=(1+g)\operatorname{id}_I,\qquad
v_1=\begin{pmatrix}1&0\\g&1\end{pmatrix},\quad
v_0=\begin{pmatrix}1&0\\1&1\end{pmatrix},\qquad
q=t\operatorname{id}_G.
\]
The relation $\operatorname{diag}(g,1)v_1
=v_0\operatorname{diag}(g,1)$ verifies that $v$ is a chain map.
All three components are isomorphisms. The four naturality equations
can be checked without the block convention:
\begin{align*}
v_1\binom11&=\binom{1}{1+g}
             =\binom t1(1+g),&
v_0\binom g1&=\binom{g}{1+g}
             =\binom{gt}{1}(1+g),\\
\begin{pmatrix}1&-t\end{pmatrix}v_1
 &=\begin{pmatrix}t&-t\end{pmatrix},&
\begin{pmatrix}1&-gt\end{pmatrix}v_0
 &=\begin{pmatrix}t&-gt\end{pmatrix}.
\end{align*}
In the paper's block notation $v$ is
$\begin{pmatrix}1&0\\1&1\end{pmatrix}$, so every component is
lower triangular. This is therefore a morphism in
$i_{\mathrm{lower}}S_2^{\triangle}\mathcal C$.

The assignments $\mu'_1,\mu'_2$ preceding Definition~12 take an
object $(n,m)_B$ to $\operatorname{Typ}(g)^{\oplus n}$ or
$\operatorname{Typ}(1)^{\oplus n}$, and a morphism to its upper-left
$n'\times n$ block in both degrees. Each $\beta_b$ consequently
maps to the identity of $\operatorname{Typ}(g)$ or
$\operatorname{Typ}(1)$. The component $v$ maps to the identity,
whereas $q$ maps to multiplication by $t$. Naturality of the
resulting right-hand square would require $1=t$. This is false,
since $1-t=g/(1+g)\ne0$. The assignments thus fail to send this
morphism of $S_2$-objects to a morphism, contrary to
\cite[Lemma~13(III)]{MochizukiParameters}.

The failure already occurs with $B=\Z_{(3)}$, $g=3$, and $t=1/4$.
The calculation uses a regular DVR with regular quotient. It identifies
a failure of the stated simplicial-functor construction used in the
subsequent additivity argument.


